\documentclass[
    bindingoffset=5mm,
    footnoteindent=3mm,
    hyphenation=true    
]{src/wut-thesis}

\graphicspath{{tex/img/}}
\addbibresource{bibliografia.bib}
\facultymeil 
\EngineerThesis 
\langpol 

\begin{document}
\instytut{samolotów i śmigłowców}
\kierunek{mechanika i projektowanie maszyn}
\specjalnosc{lotnictwo}
\title{
    Projekt koncepcjny\\bezzałogowego statku powietrznego
}

\engtitle{
    Conceptual project of\\an ultralight aircraft
}

\poltitle{
    Projekt koncepcjny bezzałogowego statku powietrznego
}
\author{Emil Guliński}
\album{316102}
\promotor{prof. dr hab. inż. Cezary Galiński}
\date{\the\year}
\maketitle

\cleardoublepage
\abstract

\indent
Niniejsza praca przedstawia wstępny projekt bezzałogowego statku powietrznego przeznaczonego do misji wywiadowczych, 
obserwacyjnych i rozpoznawczych.
Dron charakteryzuje się konfiguracją stałopłata z napędem splinowym w układzie górnopłata.
Co ułatwia prowadzenie operacji w ograniczonej przestrzeni i skraca przygotowanie do następnej operacji. 

\indent
Kluczowe elementy drona obejmują komorę ładunkową zdolną do pomieszczenia różnych czujników
takich jak kamera elektrooptyczna oraz radar SAR. 
Dron jest również wyposażony w autonomiczne systemy nawigacji oraz pokładowe jednostki przetwarzające, 
umożliwiające analizę danych i podejmowanie decyzji w czasie rzeczywistym.

\indent
W procesie projektowyn wykorzystuję następujące technologie.
Pythona3 z zaimplemontowanymi bibiotekami: 
NumPy, 
SciPy, 
SymPy, 
Matplotlib, 
Pandas do zatomatyzowania obilczeń inżynierskich oraz tworzenia wykresów.
Programu FreeCAD do stworzenia rysunku oraz modelu projektu.
W celu optymalizacji aerodynamicznej wykorzystuję komputerową mechnikę płynów (CFD) wykonaną programie OpenFOAM.
W ten sposób walidacjąc stabilności i osiągi statku. 
Konstrukcja drona została zaprojektowana tak, 
aby pomieścić niezbędne systemy przy jednoczesnym zachowaniu wydajności aerodynamicznej. 
Budowa strukturalna opiera się na kadłubie samonośnym.
Przy wykorzystaniu materiałów kompozytowych, 
co pozwala na optymalizację stosunku wytrzymałości do masy.

\keywords FreeCAD, Kompozyt, Matplotlib, NumPy, OpenFOAM, Python3, SciPy, SymPy

\clearpage
\secondabstract
\indent
This thesis presents a preliminary design of an unmanned aerial vehicle intended for intelligence missions, 
observation and reconnaissance systems.
The drone is characterized by a fixed-wing configuration with gas propulsion in a high-wing system.
This makes it easier to conduct operations in a confined space and shortens the preparation for the next operation.

\indent
The key components of the drone include a cargo compartment capable of housing various sensors
such as electro-optical camera and SAR radar. 
The drone is also equipped with autonomous navigation systems and on-board processing units, 
enabling data analysis and real-time decision-making.

\indent
I use the following technologies in the design process.
Python3 with pre-installed libraries: 
NumPy, 
SciPy, 
SymPy, 
Matplotlib, 
Pandas for automatizing engineering and creating graphs.
FreeCAD to create a drawing and a model of the project.
For aerodynamic optimization, I use computer fluid mechanics (CFD) made with OpenFOAM.
Thus, validating the stability and performance of the ship. 
The design of the drone is designed in such a way that 
to accommodate the necessary systems while maintaining aerodynamic efficiency. 
The structural structure is based on a self-supporting hull.
With the use of composite materials, 
which allows you to optimize the strength-to-weight ratio.
\secondkeywords FreeCAD, Composite, Matplotlib, NumPy, OpenFOAM, Python3, SciPy, SymPy

\cleardoublepage  
\pagestyle{plain}
\makeauthorship

\cleardoublepage 
\tableofcontents

\cleardoublepage 
\pagestyle{headings}
\section{Wstęp}

    \paragraph{}
        Bezzałogowe statki powietrzne stają się jednym z kluczowych elementów współczesnych systemów obserwacyjnych i bojowych.
        Ich zdolność do wielogodzinnego lotu nad trudno dostępnym terenem sprawia,
        że znajdują zastosowanie zarówno w misjach cywilnych,
        jak i operacjach militarnych.

    \paragraph{}
        Rosnące znaczenie tych platform potwierdzają dwa współczesne konflikty zbrojne.
        W wojnie rosyjsko-ukraińskiej drony szczególności FPV z narzędzia rozpoznawczego przekształciły się w niedozowny element uzbrojenia.
        Zyskując na statusie jak legendarny karabinek AK konstrukcji Michaiła Kałasznikowa
        Tworząc rozległe strefy kontroli przestrzeni powietrznej wzdłuż całej linii frontu.
        Ukraina zbudowała wokół nich odrębny rodzaj oddziałów wojskuch np USF Unmanned Systems Forces.
        Z kolei konflikt w Tigraju Etiopia, 2020 2022 pokazuje,
        że dostęp do platform BSP może dosłownie odwrócić losy kampanii: 
        rząd etiopski, 
        stojący w obliczu militarnej klęski.
        Pozyskał chińskie drony Wing Loong 2 oraz tureckie Bayraktar TB2, 
        które zatrzymały natarcie sił Tigraju i umożliwiły kontrofensywę.

    \paragraph{}
        Oba przypadki dowodzą tej samej zasady strona dysponująca stałym rozpoznaniem 
        z powietrza i możliwością precyzyjnego rażenia uzyskuje przewagę operacyjną.
        Której nie równoważy liczebna wyższość wojsk naziemnych, 
        które zreszta wukorzystuja drony tylko tew małe.
        W niniejszej pracy proces obliczeń został zrealizowany przy użyciu języka Python, który pozwolił do bardxziej precyzujnych obliczeń.
        Celem pracy jest opracowanie bezzałogowego systemu klasy MALE o wysokiej autonomii lotu,
        modułowej konstrukcji i możliwości integracji z różnorodnymi sensorami, 
        przeznaczonego do misji rozpoznawczych, 
        patrolowania granic oraz wsparcia działań taktycznych.

    \paragraph{}
        Conflicts such as the ISIS insurgency, the Indo-Pak conflict over Kashmir, Russian
        military actions in Ukraine, the Syrian civil war, the Lebanon conflict, and the tensions between
        the US and North Korea have heightened security concerns. These factors contribute positively
        to the market, thus driving its growth during the forecast period
\clearpage
\section{Analiza trendów}

    \subsection{Wykresy predykcjne}

    \paragraph{}
            W celu przeprowadzenia analizy zgromadziłem dane techniczne 27 dronów klasy MALE. 
            Na ich podstawie opracuję wykresy kluczowych parametrów. 
            Pierwszym etapem będzie wyznaczenie zależności czasu lotu (jako najważniejszego parametru tej klasy) od roku produkcji płatowca. 
            W dalszej kolejności przeanalizuję wpływ wydłużenia skrzydła na najważniejsze osiągi, 
            obciążenie powierzchni oraz ciąg, 
            a także zbadam korelacje między innymi parametrami projektowymi. 
            Szczegółowe dane wejściowe wykorzystane do wizualizacji znajdują się w Załączniku nr 1.

        \begin{figure}[!h]
            \centering \includegraphics[width=1\linewidth]{Wykresy/Czas_lotu_w_zależności_od_roku_produkcji.png}
            \caption{Wykres zależności czasu od roku produkcji}
        \end{figure}

        \begin{figure}[!h]
            \centering \includegraphics[width=1\linewidth]{Wykresy/Czas_lotu_w_zależności_od_obciążenia_mocy.png}
            \caption{Wykres zależności obciążenia mocy od czasu lotu}
        \end{figure}

        \begin{figure}[!h]
            \centering \includegraphics[width=1\linewidth]{Wykresy/Czas_lotu_w_zależności_od_obciążenia_powierzchni.png}
            \caption{Wykres zależności obciążenia powierzchni od czasu lotu}
        \end{figure}

        \begin{figure}[!h]
            \centering \includegraphics[width=1\linewidth]{Wykresy/Czas_lotu_w_zależności_od_wydłużenia.png}
            \caption{Wykres zależności obciążenia powierzchni od czasu lotu}
        \end{figure}

        \begin{figure}[!h]
            \centering \includegraphics[width=1\linewidth]{Wykresy/Masa_własna_startowa_w_zależności_od_masy_startowej.png}
            \caption{Wykres zależności wspołczunika masy własnej do startowej w zależnosći od masy startowej}
        \end{figure}
    
    \newpage
    \subsection{Podsumowanie predykcji}

        \paragraph{}
            Wykresy pozwoliły mi na oszacowanie minimalnej wartości najważniejszego parametru statku powietrznego
            czasu lotu a także pozostałych współczynników opisujących charakterystykę dynamiczną BSP. 
            Na podstawie wykresu przedstawiającego stosunek masy własnej do masy startowej w funkcji masy startowej wyznaczyłem współczynniki $A$ i $C$, 
            szerzej opisane w literaturze \cite{galinski}. 
            Poprawność wyznaczonych współczynników zweryfikowałem następnie poprzez porównanie ich 
            z wartościami dostępnymi w literaturze przedmiotu \cite{raymer}.

                  \begin{longtable}{| c | m{0,5\linewidth} | c |}
            \caption{Współczyniki opisujące charakterystykę dynamiczną}
            \label{table:współczynników opisujących charakterystykę dynamiczną} \\
            \hline
            \textbf{Nazwa parametru} & \multicolumn{1}{c|}{\textbf{Wartość}} & \textbf{Jednostka} \\ \hline \hline
            \endfirsthead
            
            \hline
            \endfoot
            
            Czas lotu & 42,135 & h \\ \hline
            Wydłużenie pałata & 15,468 & N/A \\ \hline
            Obciązenie powierzchni & 97,177 & kg/m\textsuperscript{2} \\ \hline
            Obciązenie mocy & 11,648 & kg/kW \\ \hline
            Współczynik A & 1,510 & N/A \\ \hline
            Współczynik C & -0,135 & N/A \\ \hline

        \end{longtable}
\clearpage
\section{Analiza konkurecjnych konstrukcji}
    \paragraph{}
        Przed wykonaniem odrecznego rysunku prototypo wyselekcjonowałem 5 najbardziej zblizonych konstrukcje.
        Selekcji dokonałem za pomoca odległosci euklidesowej opisanej w literaturze \cite{analnaIza}.
        
    \subsection{Milkor 380}

        \begin{figure}[!h]
            \centering \includegraphics[width=1\linewidth]{Drony/Dron1/1.jpg}
            \caption{Milkor 380\cite{Dron1}}
        \end{figure}

        \begin{longtable}{| c | m{0,5\linewidth} | c |}
            \caption{Szczegółowe parametry Milkor 380}
            \label{table:milkor_380_specs} \\
            \hline
            \textbf{Lp} & \multicolumn{1}{c|}{\textbf{Nazwa parametru}} & \textbf{Wartość} \\ \hline \hline
            \endfirsthead
            
            \hline
            \endfoot
            
            1 & Producent & Milkor \\ \hline
            2 & Kraj produkcji & RPA \\ \hline
            3 & Dopasowanie & 93,23\% \\ \hline
            4 & Model silnika & Rotax 915 iS \\ \hline
            6 & Zasięg maksymalny  & 4000,0 km \\ \hline
            7 & Prędkość przelotowa / maksymalna & 150 / 250 km/h \\ \hline
            8 & Pułap operacyjny & 9144 m \\ \hline
            9 & Zasięg operacyjny SATA & 4000 km \\ \hline
            10 & Zasięg operacyjny LOS & 220 km \\ \hline
            11 & Długość / Wysokość / Rozpiętość & 9,00 / 2,40 / 18,60 m \\ \hline
            15 & Profil skrzydła root / tip & NACA 64-415 / NACA 64-412 \\ \hline
            16 & Masa własna & 700,0 kg \\ \hline
            17 & Masa startowa & 1300,0 kg \\ \hline
            18 & Maksymalna ładowność & 210,0 kg \\ \hline

        \end{longtable}
    
    \clearpage
    \subsection{Elbit Hermes 900}
        \begin{figure}[!h]
            \centering \includegraphics[width=1\linewidth]{Drony/Dron2/1.jpg}
            \caption{Elbit Hermes 900\cite{Dron2}}
        \end{figure}

        \begin{longtable}{| c | m{0,5\linewidth} | c |}
            \caption{Szczegółowe parametry Elbit Hermes 900}
            \label{table:hermes_900_specs} \\
            \hline
            \textbf{Lp} & \multicolumn{1}{c|}{\textbf{Nazwa parametru}} & \textbf{Wartość} \\ \hline \hline
            \endfirsthead
            
            \hline
            \endfoot
            
            1 & Producent & Elbit Systems \\ \hline
            2 & Kraj produkcji & Izrael \\ \hline
            3 & Dopasowanie & 89,99\% \\ \hline
            4 & Model silnika & Rotax 914 \\ \hline
            5 & Czas lotu & 36,0 h \\ \hline
            6 & Zasięg maksymalny & 5000,0 km \\ \hline
            7 & Prędkość przelotowa / maksymalna & 112 / 220 km/h \\ \hline
            8 & Pułap operacyjny & 9144 m \\ \hline
            9 & Zasięg operacyjny SATA & 5000 km \\ \hline
            10 & Zasięg operacyjny LOS & 250 km \\ \hline
            11 & Długość / Wysokość / Rozpiętość & 8,30 / 2,30 / 15,00 m \\ \hline
            12 & Wydłużenie płata & 15,00 \\ \hline
            13 & Obciążenie powierzchni & 78,67 kg/m$^2$ \\ \hline
            14 & Obciążenie mocy & 13,95 kg/kW \\ \hline
            15 & Profil skrzydła root / tip & Wortmann FX 63-137 \\ \hline
            16 & Masa własna & 750,0 kg \\ \hline
            17 & Masa startowa & 1180,0 kg \\ \hline
            18 & Maksymalna ładowność & 350,0 kg \\ \hline
        \end{longtable}

    \clearpage
    \subsection{Baykar Bayraktar TB2}

        \begin{figure}[!h]
            \centering \includegraphics[width=1\linewidth]{Drony/Dron3/1.jpg}
            \caption{Baykar Bayraktar TB2\cite{Dron3}}
        \end{figure}

        \begin{longtable}{| c | m{0,5\linewidth} | c |}
            \caption{Szczegółowe parametry Baykar Bayraktar TB2}
            \label{table:bayraktar_tb2_specs} \\
            \hline
            \textbf{Lp} & \multicolumn{1}{c|}{\textbf{Nazwa parametru}} & \textbf{Wartość} \\ \hline \hline
            \endfirsthead
            
            \hline
            \endfoot
            
            1 & Producent & Baykar Teknoloji \\ \hline
            2 & Kraj produkcji & Turcja \\ \hline
            3 & Dopasowanie & 89,35\% \\ \hline
            4 & Model silnika & Rotax 912 \\ \hline
            5 & Czas lotu & 27,0 h \\ \hline
            6 & Zasięg maksymalny & 4000,0 km \\ \hline
            7 & Prędkość przelotowa / maksymalna & 130 / 220 km/h \\ \hline
            8 & Pułap operacyjny & 8229 m \\ \hline
            9 & Zasięg operacyjny SATA & 4000 km \\ \hline
            10 & Zasięg operacyjny LOS & 300 km \\ \hline
            11 & Długość / Wysokość / Rozpiętość & 6,50 / 2,20 / 12,00 m \\ \hline
            12 & Wydłużenie płata & 13,71 \\ \hline
            13 & Obciążenie powierzchni & 66,67 kg/m$^2$ \\ \hline
            14 & Obciążenie mocy & 9,52 kg/kW \\ \hline
            15 & Profil skrzydła root / tip & NACA 4412 \\ \hline
            16 & Masa własna & 420,0 kg \\ \hline
            17 & Masa startowa & 700,0 kg \\ \hline
            18 & Maksymalna ładowność & 150,0 kg \\ \hline
        \end{longtable}

    \clearpage
    \subsection{Shahed 129}

        \begin{figure}[!h]
            \centering \includegraphics[width=1\linewidth]{Drony/Dron4/1.png}
            \caption{Shahed 129\cite{Dron4}}
        \end{figure}

        \begin{longtable}{| c | m{0,5\linewidth} | c |}
            \caption{SzczegółoAwe parametry Shahed 129}
            \label{table:Shahed 129} \\
            \hline
            \textbf{Lp} & \multicolumn{1}{c|}{\textbf{Nazwa parametru}} & \textbf{Wartość} \\ \hline \hline
            \endfirsthead
            
            \hline
            \endfoot
            1 & Producent & IAMIC \\ \hline
            2 & Kraj produkcji & Iran \\ \hline
            3 & Dopasowanie & 88,30\% \\ \hline
            4 & Model silnika & Rotax 914 \\ \hline
            5 & Czas lotu & 24,0 h \\ \hline
            6 & Zasięg maksymalny & 1700,0 km \\ \hline
            7 & Prędkość przelotowa / maksymalna & 150 / 175 km/h \\ \hline
            8 & Pułap operacyjny & 7300 m \\ \hline
            9 & Zasięg operacyjny SATA & 1700 km \\ \hline
            10 & Zasięg operacyjny LOS & 200 km \\ \hline
            11 & Długość / Wysokość / Rozpiętość & 8,00 / 3,10 / 16,00 m \\ \hline
            12 & Wydłużenie płata & 14,88 \\ \hline
            13 & Obciążenie powierzchni & 63,95 kg/m$^2$ \\ \hline
            14 & Obciążenie mocy & 13,01 kg/kW \\ \hline
            15 & Masa własna & 700,0 kg \\ \hline
            16 & Masa startowa & 1100,0 kg \\ \hline
            17 & Maksymalna ładowność & 400,0 kg \\ \hline
        \end{longtable}

    \clearpage
    \subsection{Wing Loong I}

        \begin{figure}[!h]
            \centering \includegraphics[width=1\linewidth]{Drony/Dron5/1.jpg}
            \caption{Shahed 129\cite{Dron5}}
        \end{figure}

        \begin{longtable}{| c | m{0,5\linewidth} | c |}
            \caption{SzczegółoAwe parametry Shahed 129}
            \label{table:Wing Loong I} \\
            \hline
            \textbf{Lp} & \multicolumn{1}{c|}{\textbf{Nazwa parametru}} & \textbf{Wartość} \\ \hline \hline
            \endfirsthead
            
            \hline
            \endfoot
            1 & Producent & CAIG \\ \hline
            2 & Kraj produkcji & Chiny \\ \hline
            3 & Dopasowanie & 87,77\% \\ \hline
            4 & Model silnika & Rotax 914 \\ \hline
            5 & Czas lotu & 20,0 h \\ \hline
            6 & Zasięg maksymalny & 4000,0 km \\ \hline
            7 & Prędkość przelotowa / maksymalna & 150 / 280 km/h \\ \hline
            8 & Pułap operacyjny & 7500 m \\ \hline
            9 & Zasięg operacyjny SATA & 4000 km \\ \hline
            10 & Zasięg operacyjny LOS & 200 km \\ \hline
            11 & Długość / Wysokość / Rozpiętość & 9,05 / 2,70 / 14,00 m \\ \hline
            12 & Wydłużenie płata & 12,17 \\ \hline
            13 & Obciążenie powierzchni & 68,32 kg/m$^2$ \\ \hline
            14 & Obciążenie mocy & 13,01 kg/kW \\ \hline
            15 & Profil skrzydła root / tip & Drela GW-19 / Drela GW-27 \\ \hline
            16 & Masa własna & 600,0 kg \\ \hline
            17 & Masa startowa & 1100,0 kg \\ \hline
            18 & Maksymalna ładowność & 200,0 kg \\ \hline
        \end{longtable}

    \clearpage
    \subsection{Model prototypu}

    \begin{figure}[!h]
        \centering \includegraphics[width=1\linewidth]{Model/model.pdf}
        \caption{Model prototypu}
    \end{figure}

    \paragraph{}
        W strukturze projektowanego BSP zdecydowano się na zastosowanie układu centropłata, 
        Wybór ten podyktowany jest przede wszystkim pozwoli osiągnąć optymalne osiągi aerodynamiczne oraz wysokią manewrowości
        oraz wyposazony w chowane podwozie o budowie,  
        Konfiguracja ta wiąże się z ograniczeniem przestrzeni użytecznej wewnątrz kadłuba,
        Omawiany dron nie jest przeznaczony do przewozu wielkogabarytowego ładunku płatnego, 
        lecz do obserwacji, 
        Wyposażenie stanowi jedynie zintegrowana głowica optoelektroniczna,
    
    \paragraph{}
        Niemniej jednak, mając na uwadze konieczność optymalizacji czasu trwania procesu projektowego oraz dążenie do minimalizacji ryzyka błędu, 
        zdecydowano się na implementację usterzenia klasycznego, 
    
    \paragraph{}
        W projekcie zastosowano układ ze śmigłem pchającym z tyłu kadłuba, 
        Decyzja ta wynika z konieczności zapewnienia niezakłóconego pola widzenia dla głowicy optoelektronicznej, 
        Przeniesienie napędu na ogon uwalnia przestrzeń na dziobie drona, 
        eliminując ryzyko wchodzenia łopat śmigła w kadr i powstawania zakłóceń obrazu, 
        Dodatkowo chroni to delikatną optykę przed zanieczyszczeniami i uszkodzeniami podczas lądowania,
\newpage
\section{Definicjia profilu misji}

    \begin{figure}[!h]
        \centering \includegraphics[width=1\linewidth]{ProfilMisji/misja.png}
        \caption{Profil misji BSP}
    \end{figure}

    \begin{itemize}
        \item Start
        \item Wznoszenie do wyskosści przelotowej
            \begin{itemize}
                \item prędkosci wznoszenia przelotowego: 28 m/s 
                \item kat wznoszenia: 15\textdegree
            \end{itemize}
        \item Dolot
            \begin{itemize}
                \item prędkosci przelotowa: 42 m/s 
                \item wysokość prtzelotowa: 3000 m
            \end{itemize}
        \item Wznoszenie do wyskosści patrolowej 
            \begin{itemize}
                \item prędkosci wznoszenia przelotowego: 28 m/s 
                \item kat wznoszenia: 10\textdegree
            \end{itemize}
        \item Patrolowanie
            \begin{itemize}
                \item prędkosci patrolowa: 21 m/s
                \item czas patrolu: 80000 s
                \item wysokość patrolowa zalezna od misji wykrycie/rozpoznanie/identyfikacja: 13000/5200/2600 m
            \end{itemize}
        \item Opadanie
            \begin{itemize}
                \item prędkosci wznoszenia przelotowego: 28 m/s 
                \item kat opadania: 5\textdegree
            \end{itemize}
        \item Powrót
            \begin{itemize}
                \item prędkosci przelotowa: 42 m/s 
                \item wysokość prtzelotowa: 3000 m
            \end{itemize}
        \item Oczekiwanie: 
            \begin{itemize}
                \item prędkosci patrolowa: 21 m/s
                \item czas patrolu: 900 s
                \item wysokość patrolowa zalezna od misji wykrycie/rozpoznanie/identyfikacja: 13000/5200/2600 m
            \end{itemize}
        \item Opadanie:
            \begin{itemize}
                \item prędkosci wznoszenia przelotowego: 21 m/s 
                \item kat opadania: 2\textdegree
            \end{itemize}
        \item Ladowanie
    \end{itemize}
% \clearpage
\section{Model prototypu}

    \begin{figure}[!h]
        \centering \includegraphics[width=1\linewidth]{Model/model.pdf}
        \caption{Model prototypu}
    \end{figure}

    \paragraph{}
        W strukturze projektowanego BSP zdecydowano się na zastosowanie układu centropłata, 
        Wybór ten podyktowany jest przede wszystkim pozwoli osiągnąć optymalne osiągi aerodynamiczne oraz wysokią manewrowości
        oraz wyposazony w chowane podwozie o budowie,  
        Konfiguracja ta wiąże się z ograniczeniem przestrzeni użytecznej wewnątrz kadłuba,
        Omawiany dron nie jest przeznaczony do przewozu wielkogabarytowego ładunku płatnego, 
        lecz do obserwacji, 
        Wyposażenie stanowi jedynie zintegrowana głowica optoelektroniczna,
    
    \paragraph{}
        Niemniej jednak, mając na uwadze konieczność optymalizacji czasu trwania procesu projektowego oraz dążenie do minimalizacji ryzyka błędu, 
        zdecydowano się na implementację usterzenia klasycznego, 
    
    \paragraph{}
        W projekcie zastosowano układ ze śmigłem pchającym z tyłu kadłuba, 
        Decyzja ta wynika z konieczności zapewnienia niezakłóconego pola widzenia dla głowicy optoelektronicznej, 
        Przeniesienie napędu na ogon uwalnia przestrzeń na dziobie drona, 
        eliminując ryzyko wchodzenia łopat śmigła w kadr i powstawania zakłóceń obrazu, 
        Dodatkowo chroni to delikatną optykę przed zanieczyszczeniami i uszkodzeniami podczas lądowania,
% \clearpage
\section{Szacunkowa biegunowa analityczna}
    
    \paragraph{}
        W tym rodziale obliczę współczynnik oporu samolotu przy 0 współczynnik oporu samolotu,
        Wykonam to za poprzez przekształcenie wzoru współczynnik oporu samolotu,
        Odczytanie niezbędnych wartości z wcześniejszej analizy trendu i odczytachy wartości z tabel,
        Wszystkie niezbędene tabel i oraz wzory mozna znalez w \cite{galinski}, 
        Wzor po przekształceniu zamieszczam poniżej:

        \begin{align}
            C_{x0} = \frac{\pi \cdot AR \cdot e}{4 \cdot K_{\max}^2}
        \end{align}

    \paragraph{}
        Dla moich wczesniej wspomnianych parametrów wartiość biegunowej wyszła: 0,0098
% \clearpage
\section{Wyznaczenia obciążenia powierzchni i obciążenia ciągu przyszłego samolotu}

    \subsection{Wykres ograniczeń}
        \paragraph{}
            Zgodnie z rodziałem 2.5 ksiąrzki \cite{galinski}. 
            Wykonałem wykres zależności odwrotności obciążenia ciągu od obciążenia powierzchni.

            \begin{figure}[!h]
                \centering \includegraphics[width=1\linewidth]{ObciazeniePowierzchni/ObciązeniePowierzchni.png}
                \caption{Wykres zależności czasu od roku produkcji}
            \end{figure}

        \paragraph{}
            Z wykresu jestem wstanie odzczytać iż dron by osiągnąć optymalny zasięg. 
            Musi cechować się natstępującytmi parametrami:
         
            \begin{align}
                \frac{W}{S} = 671
            \end{align}

            \begin{align}
                \frac{T}{W} = 0.319
            \end{align}

        \paragraph{} 
            Na podstawię odczytanych danych stwierdzam.
            Dron by osiądnąć optymalny zasięg potrzebuje powiechnie nośna powyżej 13m^2.
            Następnnie z wzoru 2.80 z \cite{galinski} obliczyłem,
            Układ napędowy powerplant powinnien cechowac się ciągiem powyżej 2757N.
            Silnik użyty do projektu powinien cechować sie mocą powyżej 145kW.
% \clearpage
\section{Podsumowanie analiza kosztów}

    \paragraph{}
        Zgodnie z rodziałem 2.7 Analiza kosztów ksiąrzki \cite{galinski} mojego promotora. 
        Wykonałem wykres zależności odwrotności obciążenia ciągu od obciążenia powierzchni.
        napoczątku przed wykonaniem obliczen sprawdziłem ile czy wykonałem poprawnie obliczenia.
        Szacujac koszty dla zakonczonej juz produkcji 737-300 z lat powstania raportu na podstawie ktorego.
        Postała methoda wyszło mi rok 1970: 20980423 USD oraz w roku 1986: 16368673 USD co wydaje się dos prawdopodobne.
        Ponieważ smaloty były sprzedawane w tamtym okresie za 40 - 46,5 mln USD. Przyjmuje średnią cenne 43,5 mln USD i daje 2.66 spredu.

        \begin{itemize}
            \item Rok 1970:
                \begin{itemize}
                    \item Wsparcie prac badawczo rozwojowych: 381526 USD
                    \item Cena silnika i awioniki: 9041976 USD
                    \item Materiały: 16095760 USD
                    \item Próby w locie: 23137136 USD
                    \item Koszt inzynierów po przliczeniu stawek za roboczo godziny: 914704 USD
                    \item Koszt narzędzi po przliczeniu stawek za roboczo godziny: 31959501 USD
                    \item Koszt wytworzenia po przliczeniu stawek za roboczo godziny: 2413456 USD
                    \item Koszt kontroli jakości po przliczeniu stawek za roboczo godziny: 3612820 USD
                \end{itemize}
            \item Rok 1986:
                \begin{itemize}
                    \item Wsparcie prac badawczo rozwojowych: 2718297 USD
                    \item Cena silnika i awioniki: 50976528 USD
                    \item Materiały: 57448933 USD
                    \item Próby w locie: 623756360 USD
                    \item Koszt inzynierów po przliczeniu stawek za roboczo godziny: 23053374 USD
                    \item Koszt narzędzi po przliczeniu stawek za roboczo godziny: 192492401 USD
                    \item Koszt wytworzenia po przliczeniu stawek za roboczo godziny: 20577415 USD
                    \item Koszt kontroli jakości po przliczeniu stawek za roboczo godziny: 20319358 USD
                \end{itemize}
        \end{itemize}

    \subsection{Podsumowanie analizy kosztów}

        \paragraph{}
            Podejrzewam, że autor \cite{galinski} nie bez powodu, wybrał datę 1970 i 1986.
            Chciał podkreślić to, iż o ile użycie oprogramowania CAD CAM CAE.
            Straca czas i zmniejsza koszty postania produktu.
        
        \paragraph{}
            W 1970 - używanie tych sytemów było bardzo drogie,
            więc większość używała desek reślarskich.
            Kadra inżynierska teżnie przyzwyczajona.
            W latach 70. nie było komputerów osobistych. 
            Systemy CAD działały na tzw. mainframe’ach (komputerach centralnych).
            Inżynierowie pracowali na terminalach graficznych, 
            które były niesamowicie drogie – jedno stanowisko mogło kosztować tyle, 
            co luksusowy dom.
        
        \paragraph{}
            W 1986 wtedy zaczęto używać tych narzędzi "dobrze",
            czyli w sposób zintegrowany.
            Modelowanie parametryczne (1987/1988):
            Firma PTC wypuściła program Pro/ENGINEER (dziś Creo).
            To był absolutny "game changer".
            Wprowadził on logikę,
            w której zmiana jednego wymiaru aktualizowała cały model

        \paragraph{}
            Zdecydowałem się sumowac kwotę z roku 1986 nastepnie przmnożyć ją przez wskażnik CPI.
            Zakładajac, iz roczna inflacja w okresie 2026-2030 2.5 \%. Szacowany kozt jednostkowy,
            przy zerowej marży przy założeniach sprzedarzy 1000 i produkcji 2 samolotów w miesiącu.
            Wynosi 3056193 USD.

        \paragraph{}
            Pytanie do promoamatora:
            \newline
            Czy methoda wynika z ramortu z lat 90 i skupupia sie głownie na produkcji smaolotow z durali.
            Czy do procesu produkcjnego drona też sie sprawdzi kwoty wydaja sie dośc wysokie.







% \input{tex/9-wyznaczenie-masy-startowej}
% \clearpage
\section{Studium wykonalności bezzałogowego samolotu eksperymentalnego}
    \paragraph{}
        W nawiązaniu do metodologii projektowania samolotów eksperymentalnych przedstawionej w literaturze, 
        kolejnym etapem pracy jest przeprowadzenie obliczeń skróconych dla modelu badawczego w zmniejszonej skali. 
        Zgodnie z opisanymi zasadami, w celu ograniczenia kosztów oraz umożliwienia zbadania wybranej konfiguracji aerodynamicznej, 
        zdecydowano się na zaprojektowanie bezzałogowego modelu do badań w locie.
    \paragraph{}
        Do przeliczenia właściwości aerodynamicznych i masowych pomiędzy obiektem w pełnej skali 
        a modelem badawczym zastosowana zostanie teoria podobieństwa dynamicznego, oparta na skalowaniu Froude’a 
        (zgodnie z wytycznymi z tabeli 2.12). 
        Wybór tego kryterium podobieństwa pozwoli na zachowanie odpowiednich relacji sił bezwładności do sił ciężkości, 
        co jest kluczowe podczas badań w locie.
    \paragraph{}
        W dalszej części rozdziału przeprowadzona zostanie analiza doboru masy,
        obciążenia powierzchni oraz napędu dla modelowej skali, 
        z uwzględnieniem ograniczeń wynikających zarówno z budżetu, 
        jak i planowanych parametrów geometrycznych (dążenie do uzyskania jak największej liczby Reynoldsa).

%---------------
% Bibliografia
%---------------
\cleardoublepage 
\printbibliography

% Wykaz symboli i skrótów.
% Pamiętaj, żeby posortować symbole alfabetycznie
% we własnym zakresie. Makro \acronymlist
% generuje właściwy tytuł sekcji, w zależności od języka.
% Makro \acronym dodaje skrót/symbol do listy,
% zapewniając podstawowe formatowanie.
\acronymlist
\acronym{BSP}{Bezałogowy statek powietrzny}
\acronym{SAR}{Synthetic Aperture Radar}
\acronym{UAS}{Unmanned Aircraft System}
\acronym{ISR}{Intelligence, Surveillance, and Reconnaissance}
\acronym{ISTAR}{Intelligence, Surveillance, Target Acquisition, and Reconnaissance (Rozszerzone o wskazywanie celów)}
\acronym{UCAV}{Unmanned Combat Aerial Vehicle}
\acronym{CAS}{Close Air Support}
\acronym{SEAD}{Suppression of Enemy Air Defense}
\acronym{EW}{Electronic Warfare}
\acronym{C2 / C4ISR}{Command and Control / Communications}


\vspace{0.8cm}

%--------------------------------------
% Spisy: rysunków, tabel, załączników
%--------------------------------------
\pagestyle{plain}

\listoffigurestoc    % Spis rysunków.
\vspace{1cm}         % vertical space
\listoftablestoc     % Spis tabel.
\vspace{1cm}         % vertical space
\listofappendicestoc % Spis załączników

%-------------
% Załączniki
%-------------

% Obrazki i tabele w załącznikach nie trafiają do spisów
\captionsetup[figure]{list=no}
\captionsetup[table]{list=no}

% Załącznik 1
\clearpage
\appendix{Dane do wykresów}
\begin{longtable}{| c | m{0,5\linewidth} | c | c |}
    \caption{Czas lotu i rok produkcji}
    \label{table:timeyear} \\
    \hline
    \textbf{Lp} & \multicolumn{1}{c|}{\textbf{Nazwa}} & \textbf{Rok kontraktu} & \textbf{Czas lotu [h]} \\ \hline \hline
    \endfirsthead


    \hline
    \endfoot
    1 & Aerostar & 2000 & 12 \\ \hline
    2 & Dominator & 2011 & 28 \\ \hline
    3 & Baykar Bayraktar TB2 & 2011 & 27 \\ \hline
    4 & Wing Loong I & 2010 & 20 \\ \hline
    5 & Wing Loong II & 2017 & 20 \\ \hline
    6 & CH-5 & 2017 & 60 \\ \hline
    7 & Wing Loong III & 2026 & 40 \\ \hline
    8 & FH-95 & 2021 & 24 \\ \hline
    9 & Elbit Hermes 900 & 2010 & 36 \\ \hline
    10 & General Atomics MQ-1 Predator & 1995 & 24 \\ \hline
    11 & HAI Pegasus II & 2005 & 15 \\ \hline
    12 & Heron TP & 2006 & 30 \\ \hline
    13 & KAL KUS-FS & 2023 & 24 \\ \hline
    14 & Orion & 2020 & 24 \\ \hline
    15 & Milkor 380 & 2024 & 35 \\ \hline
    16 & MQ-1C Gray Eagle Standard & 2005 & 25 \\ \hline
    17 & MQ-1C Gray Eagle Extended Range & 2015 & 42 \\ \hline
    18 & MQ-9 Reaper & 2006 & 27 \\ \hline
    19 & MQ-9B SkyGuardian & 2018 & 40 \\ \hline
    20 & Qods Mohajer-6 & 2017 & 12 \\ \hline
    21 & RQ-5 Hunter & 1993 & 12 \\ \hline
    22 & Shahed 129 & 2013 & 24 \\ \hline
    23 & Shahed 149 & 2022 & 35 \\ \hline
    24 & TAI Aksungur & 2021 & 50 \\ \hline
    25 & TAI Anka-S & 2013 & 24 \\ \hline
    26 & Vestel Karayel & 2014 & 20 \\ \hline
    27 & Yabhon United 40 & 2013 & 120 \\ \hline
\end{longtable}

\begin{longtable}{| c | m{0,5\linewidth} | c |}
    \caption{Obciążenie mocy [kg/kW]}
    \label{table:power_load} \\
    \hline
    \textbf{Lp} & \multicolumn{1}{c|}{\textbf{Nazwa}} & \textbf{Obciążenie mocy [kg/kW]} \\ \hline \hline
    \endfirsthead
    
    \hline
    \endfoot
    
    1 & Aerostar & 8,23 \\ \hline
    2 & Dominator & 19,24 \\ \hline
    3 & Baykar Bayraktar TB2 & 9,52 \\ \hline
    4 & Wing Loong I & 13,01 \\ \hline
    5 & Wing Loong II & 9,52 \\ \hline
    6 & CH-5 & 13,60 \\ \hline
    7 & Wing Loong III & 8,43 \\ \hline
    8 & FH-95 & 13,60 \\ \hline
    9 & Elbit Hermes 900 & 13,95 \\ \hline
    10 & General Atomics MQ-1 Predator & 12,06 \\ \hline
    11 & HAI Pegasus II & 8,94 \\ \hline
    12 & Heron TP & 6,42 \\ \hline
    13 & KAL KUS-FS & 6,51 \\ \hline
    14 & Orion & 11,33 \\ \hline
    15 & Milkor 380 & 12,53 \\ \hline
    16 & MQ-1C Gray Eagle Standard & 13,46 \\ \hline
    17 & MQ-1C Gray Eagle Extended Range & 14,39 \\ \hline
    18 & MQ-9 Reaper & 6,82 \\ \hline
    19 & MQ-9B SkyGuardian & 8,11 \\ \hline
    20 & Qods Mohajer-6 & 8,16 \\ \hline
    21 & RQ-5 Hunter & 9,40 \\ \hline
    22 & Shahed 129 & 13,01 \\ \hline
    23 & Shahed 149 & 5,62 \\ \hline
    24 & TAI Aksungur & 13,20 \\ \hline
    25 & TAI Anka-S & 12,80 \\ \hline
    26 & Vestel Karayel & 8,57 \\ \hline
    27 & Yabhon United 40 & 17,73 \\ \hline

\end{longtable}

\begin{longtable}{| c | m{0,5\linewidth} | c |}
    \caption{Obciążenie powierzchni {[}kg/m\textsuperscript{2}{]}}
    \label{table:surface_load} \\
    \hline
    \textbf{Lp} & \multicolumn{1}{c|}{\textbf{Nazwa}} & \textbf{\begin{tabular}[c]{@{}c@{}}Obciążenie powierzchni {[}kg/m\textsuperscript{2}{]}\end{tabular}} \\ \hline \hline
    \endfirsthead
    
    \hline
    \endfoot
    
    1 & Aerostar & 45,10 \\ \hline
    2 & Dominator & 117,25 \\ \hline
    3 & Baykar Bayraktar TB2 & 66,67 \\ \hline
    4 & Wing Loong I & 68,32 \\ \hline
    5 & Wing Loong II & 158,49 \\ \hline
    6 & CH-5 & 103,13 \\ \hline
    7 & Wing Loong III & 147,62 \\ \hline
    8 & FH-95 & 95,65 \\ \hline
    9 & Elbit Hermes 900 & 78,67 \\ \hline
    10 & General Atomics MQ-1 Predator & 89,08 \\ \hline
    11 & HAI Pegasus II & 55,56 \\ \hline
    12 & Heron TP & 147,27 \\ \hline
    13 & KAL KUS-FS & 166,67 \\ \hline
    14 & Orion & 86,96 \\ \hline
    15 & Milkor 380 & 65,66 \\ \hline
    16 & MQ-1C Gray Eagle Standard & 132,76 \\ \hline
    17 & MQ-1C Gray Eagle Extended Range & 131,38 \\ \hline
    18 & MQ-9 Reaper & 414,17 \\ \hline
    19 & MQ-9B SkyGuardian & 156,63 \\ \hline
    20 & Qods Mohajer-6 & 58,82 \\ \hline
    21 & RQ-5 Hunter & 75,00 \\ \hline
    22 & Shahed 129 & 63,95 \\ \hline
    23 & Shahed 149 & 98,41 \\ \hline
    24 & TAI Aksungur & 92,18 \\ \hline
    25 & TAI Anka-S & 117,65 \\ \hline
    26 & Vestel Karayel & 53,39 \\ \hline
    27 & Yabhon United 40 & 34,09 \\ \hline
    
\end{longtable}

\begin{longtable}{| c | m{0,5\linewidth} | c |}
    \caption{Wydłużenie płata}
    \label{table:aspect_ratio} \\
    \hline
    \textbf{Lp} & \multicolumn{1}{c|}{\textbf{Nazwa}} & \textbf{\begin{tabular}[c]{@{}c@{}}Wydłużenie płata {[}-{]}\end{tabular}} \\ \hline \hline
    \endfirsthead
    
    \hline
    \endfoot
    
    1 & Aerostar & 14,84 \\ \hline
    2 & Dominator & 11,06 \\ \hline
    3 & Baykar Bayraktar TB2 & 13,71 \\ \hline
    4 & Wing Loong I & 12,17 \\ \hline
    5 & Wing Loong II & 15,86 \\ \hline
    6 & CH-5 & 13,78 \\ \hline
    7 & Wing Loong III & 13,71 \\ \hline
    8 & FH-95 & 12,52 \\ \hline
    9 & Elbit Hermes 900 & 15,00 \\ \hline
    10 & General Atomics MQ-1 Predator & 19,13 \\ \hline
    11 & HAI Pegasus II & 8,54 \\ \hline
    12 & Heron TP & 17,56 \\ \hline
    13 & KAL KUS-FS & 18,12 \\ \hline
    14 & Orion & 23,10 \\ \hline
    15 & Milkor 380 & 17,47 \\ \hline
    16 & MQ-1C Gray Eagle Standard & 23,50 \\ \hline
    17 & MQ-1C Gray Eagle Extended Range & 21,61 \\ \hline
    18 & MQ-9 Reaper & 35,13 \\ \hline
    19 & MQ-9B SkyGuardian & 15,91 \\ \hline
    20 & Qods Mohajer-6 & 9,80 \\ \hline
    21 & RQ-5 Hunter & 9,22 \\ \hline
    22 & Shahed 129 & 14,88 \\ \hline
    23 & Shahed 149 & 14,00 \\ \hline
    24 & TAI Aksungur & 16,09 \\ \hline
    25 & TAI Anka-S & 22,52 \\ \hline
    26 & Vestel Karayel & 9,34 \\ \hline
    27 & Yabhon United 40 & 9,09 \\ \hline
\end{longtable}

\begin{longtable}{| c | m{0,5\linewidth} | c |}
    \caption{Tabela parametrów systemów bezzałogowych}
    \label{table:mass_ratio} \\
    \hline
    \textbf{Lp} & \multicolumn{1}{c|}{\textbf{Nazwa}} & \textbf{\begin{tabular}[c]{@{}c@{}}Masa własna/\\ Masa startowa\end{tabular}} \\ \hline \hline
    \endfirsthead
    
    \hline
    \endfoot
    
    1 & Aerostar & 0,74 \\ \hline
    2 & Dominator & 0,63 \\ \hline
    3 & Baykar Bayraktar TB2 & 0,60 \\ \hline
    4 & Wing Loong I & 0,55 \\ \hline
    5 & Wing Loong II & 0,39 \\ \hline
    6 & CH-5 & 0,33 \\ \hline
    7 & Wing Loong III & 0,63 \\ \hline
    8 & FH-95 & 0,55 \\ \hline
    9 & Elbit Hermes 900 & 0,64 \\ \hline
    10 & General Atomics MQ-1 Predator & 0,50 \\ \hline
    11 & HAI Pegasus II & 0,74 \\ \hline
    12 & Heron TP & 0,48 \\ \hline
    13 & KAL KUS-FS & 0,61 \\ \hline
    14 & Orion & 0,60 \\ \hline
    15 & Milkor 380 & 0,54 \\ \hline
    16 & MQ-1C Gray Eagle Standard & 0,81 \\ \hline
    17 & MQ-1C Gray Eagle Extended Range & 0,47 \\ \hline
    18 & MQ-9 Reaper & 0,47 \\ \hline
    19 & MQ-9B SkyGuardian & 0,48 \\ \hline
    20 & Qods Mohajer-6 & 0,70 \\ \hline
    21 & RQ-5 Hunter & 0,61 \\ \hline
    22 & Shahed 129 & 0,64 \\ \hline
    23 & Shahed 149 & 0,45 \\ \hline
    24 & TAI Aksungur & 0,55 \\ \hline
    25 & TAI Anka-S & 0,72 \\ \hline
    26 & Vestel Karayel & 0,75 \\ \hline
    27 & Yabhon United 40 & 0,35 \\ \hline
\end{longtable}
% Załącznik 2
\clearpage
\appendix{Nazwa załącznika 2}


% Używając powyższych spisów jako szablonu,
% możesz dodać również swój własny wykaz,
% np. spis algorytmów.

\end{document} % Dobranoc.
